\section*{Capítulo \ref{ch:b2:plant-flowering} --- Desenvolvimento reprodutivo e controle da floração}

\begin{solution}{exo:b2:plant-flowering:1}
\omterm{def:b2:plant-meristems:meristem}{Meristema} vegetativo: produz folhas com gemas axilares,
indefinidamente (indeterminado). \omterm{def:b2:plant-flowering:transition}{Meristema de inflorescência}: produz
brácteas e \omterm{def:b2:plant-meristems:meristem}{meristemas} florais em suas axilas, indeterminado em muitas
espécies. \omterm{def:b2:plant-meristems:meristem}{Meristema} floral: produz os quatro verticilos e para ---
determinado.
\end{solution}

\begin{solution}{exo:b2:plant-flowering:2}
As \omterm{prop:b2:plant-flowering:photoperiod}{plantas de dia curto} florescem quando o dia é mais curto que um
comprimento crítico (crisântemo, soja, arroz); as de dia longo, quando
é mais longo (espinafre, trigo, \emph{Arabidopsis}); as neutras
florescem ao atingir certo tamanho (tomate, milho). Uma \omterm{prop:b2:plant-flowering:photoperiod}{planta de dia
curto} mede o comprimento da noite ininterrupta.
\end{solution}

\begin{solution}{exo:b2:plant-flowering:3}
Tipo selvagem: sépala, pétala, \omterm{def:b2:angiosperm-reproduction:flower}{estame}, \omterm{def:b2:angiosperm-reproduction:flower}{carpelo}. Sem A: \omterm{def:b2:angiosperm-reproduction:flower}{carpelo}, \omterm{def:b2:angiosperm-reproduction:flower}{estame},
\omterm{def:b2:angiosperm-reproduction:flower}{estame}, \omterm{def:b2:angiosperm-reproduction:flower}{carpelo}. Sem B: sépala, sépala, \omterm{def:b2:angiosperm-reproduction:flower}{carpelo}, \omterm{def:b2:angiosperm-reproduction:flower}{carpelo}. Sem C: sépala,
pétala, pétala, sépala, e a \omterm{def:b2:angiosperm-reproduction:flower}{flor} se repete no interior.
\end{solution}

\begin{solution}{exo:b2:plant-flowering:4}
Semanas de frio e umidade (o inverno passou); luz (a \omterm{def:b2:angiosperm-reproduction:seed}{semente} está perto
da superfície); temperaturas alternadas (solo nu, sem dossel);
lixiviação pela chuva (as chuvas chegaram); escarificação (passagem por
um tubo digestório ou abrasão no solo); fumaça e calor (um incêndio
limpou o terreno).
\end{solution}

\begin{solution}{exo:b2:plant-flowering:5}
16/8: noite de 8 horas, abaixo das 10 críticas: não. 12/12: sim.
12 de luz, 6 de escuro, lampejo, 6 de escuro: duas noites de 6 horas:
não. 8 de luz, 4 de escuro, 4 de luz, 8 de escuro: noite mais longa de
8 horas: não.
\end{solution}

\begin{solution}{exo:b2:plant-flowering:6}
O \omterm{def:b2:plant-flowering:florigen}{florígeno} é um sinal transmissível por enxertia, produzido na folha
e igual nas espécies de dia curto e de dia longo --- a diferença está
em quando a folha o produz, e não no que ela produz. Com o floema
interrompido, o enxerto não induz a floração: o sinal viaja pelo
floema.
\end{solution}

\begin{solution}{exo:b2:plant-flowering:7}
$n = \ln(1/0.15)/0.05 = 1.90/0.05 = 38$ dias frios. Os períodos somam
$47$: a planta fica vernalizada no 11.º dia do terceiro período.
Com $\alpha = 0.02$, precisaria de 95 dias e não ficaria vernalizada.
\end{solution}

\begin{solution}{exo:b2:plant-flowering:8}
O gene do repressor é silenciado por marcas da cromatina que são
copiadas a cada mitose, de modo que todas as células do caule se
lembram do inverno; na linhagem germinativa e no embrião jovem, as
marcas são apagadas e o gene volta a ser expresso. É útil porque uma
\omterm{def:b2:angiosperm-reproduction:seed}{semente} liberada no verão não deve herdar a permissão de florescer do
genitor: ela precisa esperar um inverno próprio.
\end{solution}

\begin{solution}{exo:b2:plant-flowering:9}
V: germina (o fitocromo é convertido na forma que absorve o vermelho
extremo, Pfr, a forma ativa). V depois VE: não (convertido de volta em
Pr). V, VE, V: germina. V, VE, V, VE: não. Só o último lampejo conta,
porque cada conversão é completa e a \omterm{def:b2:angiosperm-reproduction:seed}{semente} lê o estado final do
pigmento.
\end{solution}

\begin{solution}{exo:b2:plant-flowering:10}
$p = 0.7$: $r = 1.28$, $1.42$, $1.40$, $-\infty$ para $g = 0.4, 0.6, 0.8,
1$: o melhor fica perto de $g = 0.7$. $p = 0.95$: $1.99$, $2.33$, $2.55$,
$-\infty$: o ótimo se desloca para $g = 1$ (cerca de $0.95$), apenas uma
reserva simbólica. Os desertos, onde os anos bons são raros, favorecem
uma \omterm{def:b2:angiosperm-reproduction:seed}{dormência} intensa e um grande banco de \omterm{def:b2:angiosperm-reproduction:seed}{sementes}; os prados, onde a
maioria dos anos é boa, favorecem a germinação de quase tudo.
\end{solution}

\begin{solution}{exo:b2:plant-flowering:11}
B em todos os verticilos: pétala, pétala, \omterm{def:b2:angiosperm-reproduction:flower}{estame}, \omterm{def:b2:angiosperm-reproduction:flower}{estame}. C em todos os
verticilos (C reprime A): \omterm{def:b2:angiosperm-reproduction:flower}{carpelo}, \omterm{def:b2:angiosperm-reproduction:flower}{estame}, \omterm{def:b2:angiosperm-reproduction:flower}{estame}, \omterm{def:b2:angiosperm-reproduction:flower}{carpelo}. Sem A e sem
B: C sozinha em toda parte --- \omterm{def:b2:angiosperm-reproduction:flower}{carpelos} nos quatro verticilos. C também
encerra o \omterm{def:b2:plant-meristems:meristem}{meristema}; sem ela, o centro permanece indeterminado e forma
\omterm{def:b2:angiosperm-reproduction:flower}{flor} dentro de \omterm{def:b2:angiosperm-reproduction:flower}{flor}.
\end{solution}

\begin{solution}{exo:b2:plant-flowering:12}
O comprimento do dia é integrado ao longo das noites pela coincidência
da luz com o relógio, e a indução exige várias dessas noites; o frio é
integrado como marcas de cromatina que se acumulam; o banco de \omterm{def:b2:angiosperm-reproduction:seed}{sementes}
integra ao longo dos anos, distribuindo a germinação. Um integrador
ignora um único dia anômalo, um degelo em janeiro ou um ano ruim; um
limiar aplicado ao valor instantâneo faria a planta florescer numa
semana quente de fevereiro, ou germinar todas as \omterm{def:b2:angiosperm-reproduction:seed}{sementes} na primeira
chuva.
\end{solution}

\begin{solution}{pb:b2:plant-flowering:1}
\textbf{1.} Uma noite de 8 horas é mais curta que as 9.5 horas críticas:
não há floração. Para vender em agosto, o produtor precisa alongar as
noites artificialmente com o pano preto.
\textbf{2.} Uma noite de 12 horas. Uma abertura em 15 de agosto
significa uma indução completada 8 semanas antes, por volta de 20 de
junho; as doze noites indutivas precisam começar por volta de 8 de
junho.
\textbf{3.} Iluminar a cultura por alguns minutos no meio de cada
noite: uma noite de 14.5 horas se torna duas de cerca de 7 horas, ambas
abaixo do comprimento crítico. Um lampejo basta porque o fitocromo é
convertido em poucos minutos e a planta mede apenas o comprimento da
escuridão ininterrupta.
\textbf{4.} Um lampejo às 21:00 deixa um período escuro das
21:00 às 07:30, de 10.5 horas --- mais longo que 9.5:
esse segmento induz. A interrupção precisa ser colocada de modo que
nenhum segmento ultrapasse o comprimento crítico.
\textbf{5.} Uma noite perdida interrompe a série de doze noites
indutivas consecutivas: a contagem recomeça e a floração é adiada em
aproximadamente tantos dias quantos tinham sido acumulados. Um
carrapicho precisa de uma única noite indutiva e já estaria
comprometido.
\textbf{6.} O fitocromo não absorve a luz verde, de modo que a planta
não vê o lampejo; a luz vermelha o converte na forma ativa, que é lida
como dia; o vermelho extremo logo depois o converte de volta, de modo
que a noite não é interrompida.
\textbf{7.} A \omterm{prop:b2:plant-flowering:photoperiod}{planta de dia longo} floresce: o \omterm{def:b2:plant-flowering:florigen}{florígeno} atravessa a
união do enxerto pelo floema. Isso identifica um sinal móvel,
transmissível por enxertia, comum às duas classes de comprimento de
dia.
\textbf{8.} $n = \ln 10/0.06 = 38.4$: 39 dias frios.
\textbf{9.} 10 em novembro, 35 até o fim de janeiro, e o 39.º dia frio
cai no quarto dia frio de fevereiro: vernalizado no início de
fevereiro.
\textbf{10.} Seis dias: $f = e^{-0.36} = 0.70 > 0.10$: ele permanece
vegetativo todo o verão e não dá grãos --- o trigo de inverno precisa
ser semeado no outono.
\textbf{11.} Sem o repressor, ele floresce sem frio: pode ser semeado na
primavera e colhido no mesmo verão --- um trigo de primavera.
\textbf{12.} As marcas de silenciamento na cromatina do repressor são
copiadas a cada mitose, de modo que o ápice se lembra durante a
primavera; na \omterm{def:b2:meiosis-heredity:meiosis}{meiose} e no embrião, as marcas são apagadas, e cada
geração precisa contar o próprio inverno.
\textbf{13.} A memória precisa ficar nas células que formarão a \omterm{def:b2:angiosperm-reproduction:flower}{flor} e
que persistem durante o inverno --- o \omterm{def:b2:plant-meristems:meristem}{meristema} ---, e ela é um estado
de sua cromatina, e não um sinal móvel; as folhas que medem o
comprimento do dia caem, e exportam uma proteína em vez disso.
\textbf{14.} Classe C: sépala, pétala, pétala, sépala (A avança para o quarto
verticilo e a \omterm{def:b2:angiosperm-reproduction:flower}{flor} se repete por dentro).
\textbf{15.} Sem \omterm{def:b2:angiosperm-reproduction:flower}{estames}, portanto sem \omterm{def:b2:plant-life-cycles:heterospory}{pólen}; alguns \omterm{def:b2:angiosperm-reproduction:flower}{carpelos} se formam
onde C está fracamente ativa, de modo que algumas \omterm{def:b2:angiosperm-reproduction:seed}{sementes} se formam de
vez em quando.
\textbf{16.} Classe A: \omterm{def:b2:angiosperm-reproduction:flower}{carpelo}, \omterm{def:b2:angiosperm-reproduction:flower}{estame}, \omterm{def:b2:angiosperm-reproduction:flower}{estame}, \omterm{def:b2:angiosperm-reproduction:flower}{carpelo}.
\textbf{17.} C reprime A: o verticilo 1 se torna \omterm{def:b2:angiosperm-reproduction:flower}{carpelo}, o verticilo 2
(B e C), \omterm{def:b2:angiosperm-reproduction:flower}{estame} --- \omterm{def:b2:angiosperm-reproduction:flower}{carpelo}, \omterm{def:b2:angiosperm-reproduction:flower}{estame}, \omterm{def:b2:angiosperm-reproduction:flower}{estame}, \omterm{def:b2:angiosperm-reproduction:flower}{carpelo}, a \omterm{def:b2:angiosperm-reproduction:flower}{flor} do mutante
A.
\textbf{18.} Os genes ABC são uma família de \omterm{prop:b2:cell-differentiation:transcriptional}{fatores de transcrição}
herdada do ancestral comum de todas as plantas com \omterm{def:b2:angiosperm-reproduction:flower}{flores}; o módulo já
existia antes que as rosas, as bocas-de-leão e \emph{Arabidopsis}
divergissem, e cada uma o conservou.
\textbf{19.} B sem A nem C não especifica nenhum órgão floral: um órgão
semelhante a uma folha ou a uma bráctea.
\textbf{20.} $p = 0.8$: $r = 1.44$ ($g = 0.5$), $1.59$ ($0.7$), $1.55$
($0.9$).
\textbf{21.} Dos três, $g = 0.7$ (o verdadeiro ótimo fica perto de
$0.8$); com $g = 1$, o fator dos anos ruins é zero e a linhagem é
extinta pelo primeiro ano ruim.
\textbf{22.} $p = 0.5$: $0.51$, $0.41$, $-0.03$: $g = 0.5$ é o melhor,
e um valor ainda menor seria melhor --- mais \omterm{def:b2:angiosperm-reproduction:seed}{dormência} no deserto.
\textbf{23.} As \omterm{def:b2:angiosperm-reproduction:seed}{sementes} do produtor são armazenadas, e não enterradas,
e sobrevivem com certeza; ele não precisa diversificar riscos dentro da
lavoura, mas precisa guardar uma reserva no depósito para ressemear
após um ano perdido --- a mesma reserva, guardada por ele em vez de pelo
solo.
\textbf{24.} Uma \omterm{def:b2:angiosperm-reproduction:seed}{semente} pequena leva reservas para um caule de um ou
dois centímetros; germinando no escuro sob o solo, ela morreria antes
de alcançar a luz, por isso espera a luz --- que só a alcança perto da
superfície. A aração traz \omterm{def:b2:angiosperm-reproduction:seed}{sementes} enterradas para cima e desencadeia
uma explosão de ervas daninhas.
\textbf{25.} Noite crítica de 9.5 horas, cobertura a partir de cerca de
8 de junho para abertura em 15 de agosto; 39 dias frios; melhor
$g \approx 0.8$ para $p = 0.8$ e cerca de $0.5$ para $p = 0.5$.
\end{solution}
